\section{Appendix: Scope and Qualifications}
\label{sec:scope}

{\small
\textbf{Edition and formulary.} The study treats the Mass of the
Twenty-ninth Sunday in Ordinary Time, Year~A, as the \work{Roman Missal},
Third Edition, and the \work{Lectionary for Mass} give it for the dioceses of
the United States, in the national calendar alone, for the occurrence of
18~October 2026. Identity, occurrence and branches are set out in the
preceding appendix and recorded in \texttt{instance/manifest.md} and
\texttt{research/context.md}. Nothing in the study holds for a celebration that
replaces this formulary with another Mass.

\textbf{Text control.} The Missal formulary was read in the Latin of the
\work{Missale Romanum}, editio typica tertia (2002), in a digitally typeset
reproduction, and compared with the list of changes for the 2008 reprint in
\work{Notitiae} 44 (2008), which lists none for this week, and with the
\work{Antiphonary} excerpted from the English Missal (2010). No page of the
2008 reprint and no United States altar book was inspected, so the approved
English orations were not read in any witness, and no chant book was
collated. The Lectionary's boundaries rest on three witnesses that agree in
every figure: the national calendar for 2026, the bishops' conference's dated
page of readings, and no.~145 of the Latin \work{Ordo lectionum Missae}
(1981). The page of readings refused direct requests on 16~October 2026 and
was read in the Internet Archive's capture of it made on 15~March 2026; the
psalm's United States locator and the division of its strophes rest on that
capture, and the printed United States Lectionary was not inspected. The
\work{Nova Vulgata} was read in the Holy See's web publication, in deliveries
that are not retained; the readings it alone supports (at Isaiah 45:4,
Psalm 96 (95):5 and Philippians 2:16, and in the antiphons' comparison) are
reported as its readings. Four questions about the text remain open: how the
printed Lectionary gives the psalm response; whether the \work{Nova Vulgata}
the \work{Ordo} used read \latin{meo} or \latin{suo} at Isaiah 45:1; the
Hebrew and Greek of Isaiah 45, which were not read in the original languages;
and manuscript evidence for Matthew 22:15--21 beyond what the SBL apparatus
reports. The branch rows rest on the United States \work{General
Instruction}: chapter~II as read in its official chapter page on 16~October
2026, and chapters~VII and~VIII as read in Internet Archive captures of their
official pages made on 13 and 18~June 2026, because the publisher's site
refused direct requests for them that day. Nothing is drawn from those two
chapters beyond nos.~364--365, 372 and 374, which the branch rows cite.

\textbf{English.} All Scripture is printed in the Douay--Rheims version as
revised by Bishop Challoner, a public-domain study translation of the
Clementine Vulgate. It is not the text proclaimed in the celebration; the
controlling English is that of the \work{Lectionary for Mass for Use in the
Dioceses of the United States of America}, and for the antiphons and orations
that of the \work{Roman Missal}, Third Edition, neither of which is reproduced.
The orations and the antiphons are identified by Latin incipit and described;
the descriptions are not translations and may not be quoted as the texts'
words. Where the Lectionary appoints part of a verse, the rest is marked
beside it. English glosses of Latin passages are the editor's renderings of the
Latin printed beside them. In quotations the Latin is normalized only by
writing \textit{i} for consonantal \textit{j} and writing out the ligatures,
and the ligature of ``Caesar'' in the nineteenth-century English translations
is written out.

\textbf{Reception.} Each appointed Scripture passage has at least one ancient
and one medieval or later witness read at an exact locus; for the second
Communion antiphon the only ancient witness is Chrysostom on Matthew's
parallel. The witnesses are those named in the References, and the search
behind them is bounded and recorded in \texttt{research/scope.md}. Augustine,
Chrysostom, Justin Martyr, Cyprian, Novatian and the \work{Epistle of
Barnabas} were read in the English of the Nicene and Ante-Nicene Fathers;
Cyril of Alexandria in Payne Smith's English; Bellarmine in O'Sullivan's
abridged English of 1866, the Latin not compared. Jerome on Matthew and
Isaiah, Hilary on Matthew, Bede on Mark, and Aquinas on Matthew, Isaiah and
Philippians were read on page images of the printings named; Ambrose in an
electronic Latin text of the Migne printing; Aquinas on First Thessalonians and
on the Psalms in the electronic Latin of the Corpus Thomisticum. No witness
was collated with a modern critical edition. Not opened, among others: Origen,
Rabanus Maurus and the \work{Catena aurea} on Matthew; Tertullian; Eusebius,
Cyril of Alexandria and Theodoret on Isaiah; Theodoret and Theophylact on
First Thessalonians; Basil on Psalm 33; Cassiodorus beyond the title of
Psalm 96. Nothing is said here about who first proposed any reading, about any
reading being the Fathers' common view, or about the order in time or the
mutual dependence of the witnesses; Augustine's own remark that ``some'' read
the coin of the soul shows only that the reading was known to him. No
commentary on any of the three orations was found. Cyril of Alexandria,
Justin Martyr and Cyprian have no row in the registry of authors' standing
used for this study and are cited only in support; Cyprian, Novatian and
\work{Barnabas} are cited only as witnesses to a variant reading. The
liturgical commentaries of Schuster and Fromage, which treat this Gospel and
Psalm 17 (16):6 within a Mass of another order, were read on page images of
their printings and are not cited.

\textbf{Relations among the elements.} Four relations are stated by the
liturgical books: the first reading is chosen in relation to the Gospel, as
the \work{Ordo}'s titles show (\work{Praenotanda} 106); the psalm responds to
the first reading (\work{General Instruction} 61); the acclamation greets the
Gospel (\work{General Instruction} 62); and the second reading and the Gospel
each belong to a semi-continuous course. The Missal's antiphons and orations
serve all three years. The recurrence of the words for serving and rendering,
image and likeness, nations and rulers, knowing, seeing and sincerity across
the texts is an observation about the texts. Among the witnesses cited,
Augustine joins the saying on tribute to Psalm 4 and to 1~John 3:2, Chrysostom
and Augustine join it to Romans 13, Chrysostom joins 1~Thessalonians 1:3 to
Acts 17, Aquinas joins the saying to the first chapter of Baruch, and Bede and
Chrysostom read Mark 10:45 against the rulers of the gentiles in its own
context. Every other connection drawn between one element and another is the
editor's synthesis, grounded in a witness's statement about his own passage
where one is cited and otherwise an editorial proposal resting on the wording
of the texts. No witness cited comments on this Mass, and none joins Isaiah 45
to Matthew 22.

\textbf{The questioners.} Several Fathers read for this study pass, in
sentences near those quoted, from the men who questioned Jesus to the Jewish
people as a whole. None of those sentences is used. The study keeps to the
agents the Gospel names, the disciples of the Pharisees and the Herodians,
treats hypocrisy as a vice that belongs to anyone, and follows the Second
Vatican Council's teaching that the Jews are not to be presented as rejected or
accursed by God as if this followed from the Scriptures (\work{Nostra aetate}
4).

\textbf{The whited wall.} Augustine's definition of the whited wall as
hypocrisy (\work{Enarrationes in Psalmos} 95.3) is used as a definition only.
In its section he applies it to the Donatist party, and that application is not
used.

\textbf{Chronology.} The Date cells print the shared chronology corpus's
answers with their relations and alternatives unchanged, and the place phrases
beside them rest on the sources recorded in \texttt{research/scope.md}. The
first reading's answers date the prophet's ministry in one author's account and
the critical placement of chapters 40--55, not the oracle. The Psalter's common
composition bound is not a date for any of the three psalms; it was read in a
delivery of 16~October 2026 of the introduction it cites, not in the earlier
state the corpus cites. The responsorial psalm's superscription answer dates the
setting its Greek and Latin title names, and its referent answers date the
event the psalm was read as foretelling. The Gospel's event has no date in the
corpus. The composition answers of the Gospels and epistles are dates of the
books. Writing sites, first readers and narrative settings are kept apart.

\textbf{Rights.} The Douay--Rheims, the Clementine Vulgate, the King James
Version, Brenton's Septuagint, Swete's Septuagint, the Nicene and Post-Nicene
and Ante-Nicene Fathers, Payne Smith's Cyril, O'Sullivan's Bellarmine, the
volumes of Migne, the Venice, Parma and Li\`ege printings of Aquinas,
Wilson's and Feltoe's editions, the Harper issue of Lewis and Short's
dictionary and the \work{Catholic Encyclopedia} are in the public domain. The
Corpus Thomisticum texts are quoted in short phrases. The Latin and English of
the Missal, the \work{Ordo lectionum Missae}, the Lectionary, the
\work{General Instruction}, the national calendar, the \work{Nova Vulgata},
\work{Nostra aetate}, the World Mission Day message and the page of the
Pontifical Mission Societies, and the introductions of the \work{New American
Bible, Revised Edition} are protected; they are cited, and
only incipits, the three titles of the \work{Ordo}, its Latin response and
verse, short Latin phrases of the orations needed by the argument, and short
phrases of the Council's declaration and of the introductions are given.

\textbf{Records.} The composition of the target is audited in
\texttt{propers/verified.md}, and the Missal formulary in the collection's
record for the Twenty-ninth Week; the reception search, loci, disagreements and
negative results in \texttt{research/scope.md}; and the reasoning behind the
three interpretations, with the class of every cross-element link, in
\texttt{research/interpretations.md}. The research records received an
independent review before the study was written.
\par}

\phantomsection
\section*{References}
% A starred heading sets no mark, so without this the Appendix above would go
% on heading the pages of the bibliography. The heading stays starred because
% check-content-preflight finds the References section by that exact form.
\runninghead{References}
\addcontentsline{toc}{section}{References}
\label{sec:references}

{\small
\subsection*{Liturgical books, Bibles and texts}
\begin{itemize}[itemsep=0pt,topsep=.25em]
\item \work{Missale Romanum}, editio typica tertia (Typis Vaticanis, 2002), \latin{Dominica XXIX ``per annum''}; the rubrics of \latin{Tempus per annum} 5--6; the \latin{Normae universales de anno liturgico et de calendario} 59--60 and the \latin{Calendarium Romanum generale}, as printed in that Missal; the Mass \latin{Pro evangelizatione populorum}.
\item \work{Variationes et additiones in Missali Romano}, \work{Notitiae} 44 (2008), p.~371.
\item \work{The Roman Missal}, Third Edition, for Use in the Dioceses of the United States of America (2011), as represented by the \work{Antiphonary: Excerpted from the Roman Missal} (International Commission on English in the Liturgy, 2010), p.~81.
\item \work{Ordo lectionum Missae}, editio typica altera (Libreria Editrice Vaticana, 1981): \work{Praenotanda} 105--107; \latin{Tabella II}; no.~145, \latin{Dominica vigesima nona}, Year~A.
\item \work{Lectionary for Mass for Use in the Dioceses of the United States of America}, second typical edition, vol.~I, no.~145, as witnessed by the United States Conference of Catholic Bishops, \work{Liturgical Calendar for the Dioceses of the United States of America} 2026, pp.~5 and 41, and its page of readings for 18~October 2026, \url{https://bible.usccb.org/bible/readings/101826.cfm}, read in the Internet Archive capture of 15~March 2026.
\item \work{General Instruction of the Roman Missal}, United States edition of 2011 with the emendations of 2021, nos.~48, 51, 61--62, 74, 87 (chapter~II, the bishops' conference's official page, read on 16~October 2026), 364--365 (chapter~VII) and 372, 374 (chapter~VIII), the last two chapters read in Internet Archive captures of the official pages of 13 and 18~June 2026.
\item The Holy Bible, Douay--Rheims version, revised by Bishop Richard Challoner (Project Gutenberg text): Gn 1:26; Ps 16:1--8; 32:13--19; 57:2; 95:1--13; 145:2; Is 1:1; 44:24--45:13; 48:20; Mt 21:10, 23, 45--46; 22:15--22; 24:1; Mk 10:42--45; Acts 17:1--10; Rom 13:7; 1~Cor 15:49; Phil 1:13; 2:12--16; 4:22; 1~Thes 1:1--5; Jas 1:8.
\item \work{Biblia Sacra Vulgatae editionis} (Clementine Vulgate), eBible.org text: Ps 16:6--8; 32:18; 95:5, 7; Is 45:1--6; Mt 22:15--21; Mk 10:45; Phil 2:15--16.
\item \work{Nova Vulgata Bibliorum Sacrorum editio}, as published by the Holy See at vatican.va: Ps 17 (16):6--8; 33 (32):18; 96 (95):5, 7; Is 45:1, 4; Mk 10:45; Phil 2:16.
\item The Holy Bible, King James Version (eBible.org text): Is 45:4; Ps 96:5.
\item L.~C.~L. Brenton, \work{The Septuagint Version of the Old Testament, with an English Translation}: Is 45:4.
\item H.~B. Swete, \work{The Old Testament in Greek according to the Septuagint}: Ps 32 (title); 95:1, 5, 7.
\item \work{The Greek New Testament: SBL Edition}, ed. M.~W. Holmes, text and apparatus of Matthew at 22:15--21.
\end{itemize}

\subsection*{Fathers, Doctors and saints}
\begin{itemize}[itemsep=0pt,topsep=.25em]
\item St Augustine, \work{Enarrationes in Psalmos} 4.8; 16.1--2, 6, 8; 32.16, 18--19; 57.8; 95.1--3, 5, 8--9; 115.13; 118, under Schin; 124.7--8; Nicene and Post-Nicene Fathers, 1st ser., vol.~8.
\item St Augustine, \work{In Iohannis evangelium tractatus} 40.9 and 41.2, Nicene and Post-Nicene Fathers, 1st ser., vol.~7.
\item St Augustine, \work{De sermone Domini in monte} II.20.70, Nicene and Post-Nicene Fathers, 1st ser., vol.~6.
\item St Ambrose, \work{Expositio evangelii secundum Lucam} IX.34--36 (Migne, PL 15, cols.~1502--1503), in an electronic Latin text of that printing.
\item St Hilary of Poitiers, \work{Commentarius in Matthaeum} 23.1--2 (Migne, PL 9, Paris, 1844, cols.~1044--1045).
\item St Jerome, \work{Commentarii in Matthaeum} III, on 22:15--21 (Migne, PL 26, Paris, 1845, cols.~162--163); \work{Commentarii in Isaiam} XII, on 44:24--45:7 (Migne, PL 24, Paris, 1865, cols.~456--458).
\item St John Chrysostom, \work{Homilies on Matthew} 6, 65.4 and 70.1--2, Nicene and Post-Nicene Fathers, 1st ser., vol.~10; \work{Homilies on Romans} 23, vol.~11; \work{Homilies on Philippians} 8 and \work{Homilies on First Thessalonians} 1, vol.~13.
\item St Bede, \work{In Marci evangelium expositio} III, on 10:42--45 (Migne, PL 92, Paris, 1862, col.~236).
\item St Thomas Aquinas, \work{Super Evangelium Matthaei} c.~22, in \work{Opera}, editio altera Veneta, vol.~3 (Venice, 1745), pp.~283--284.
\item St Thomas Aquinas, \work{Expositio super Isaiam ad litteram} c.~45, in \work{Opera omnia}, vol.~14 (Parma, 1863), p.~539.
\item St Thomas Aquinas, \work{Super epistolam ad Philippenses} c.~2, lect.~4, in \work{In omnes S.~Pauli Apostoli epistolas commentaria}, vol.~2 (Li\`ege: Dessain, 1857), p.~383.
\item St Thomas Aquinas, \work{Super I epistolam ad Thessalonicenses} c.~1, lect.~1, and \work{Postilla super Psalmos} 16 and 32 (Corpus Thomisticum text).
\item St Robert Bellarmine, \work{A Commentary on the Book of Psalms}, trans. J. O'Sullivan (1866, abridged), in the eCatholic2000 transcription: on Ps 16:3, 8; 32:18--19; 95:5, 10, 13.
\end{itemize}

\subsection*{Other early witnesses, cited in support}
\begin{itemize}[itemsep=0pt,topsep=.25em]
\item Cyril of Alexandria, \work{Commentary on the Gospel of Saint Luke}, Sermon 135, trans. R. Payne Smith (Oxford, 1859).
\item St Justin Martyr, \work{First Apology} 17, Ante-Nicene Fathers, vol.~1.
\item \work{Epistle of Barnabas} 12, Ante-Nicene Fathers, vol.~1.
\item St Cyprian, \work{Ad Quirinum} (\work{Testimonia}) I.21, and Novatian, \work{De Trinitate} 26, Ante-Nicene Fathers, vol.~5.
\end{itemize}

\subsection*{Sacramentaries}
\begin{itemize}[itemsep=0pt,topsep=.25em]
\item H.~A. Wilson, ed., \work{The Gelasian Sacramentary: Liber Sacramentorum Romanae Ecclesiae} (Oxford, 1894), pp.~105 and 229, for the Collect.
\item C.~L. Feltoe, ed., \work{Sacramentarium Leonianum} (Cambridge, 1896), pp.~17 and 126, for the Prayer over the Offerings and the Prayer after Communion.
\end{itemize}

\subsection*{Magisterium}
\begin{itemize}[itemsep=0pt,topsep=.25em]
\item Second Vatican Council, Declaration on the Relation of the Church to Non-Christian Religions \work{Nostra aetate} (28~October 1965), no.~4, English and Latin texts of the Holy See's website.
\item Leo XIV, \work{Message for the 100th World Mission Day 2026 [18 October 2026]: One in Christ, united in mission} (25~January 2026), English text of the Holy See's website, read on 16~October 2026.
\end{itemize}

\subsection*{Other works}
\begin{itemize}[itemsep=0pt,topsep=.25em]
\item C.~T. Lewis and C. Short, \work{Harpers' Latin Dictionary: A New Latin Dictionary} (New York: Harper \& Brothers, copyright 1879), s.v. \latin{majestas}, p.~1101.
\item The Pontifical Mission Societies in the United States, \work{World Mission Sunday 2026}, \url{https://pontificalmissions.org/wms-2026}, read on 16~October 2026.
\end{itemize}

\subsection*{Chronology sources}
\begin{itemize}[itemsep=0pt,topsep=.25em]
\item \work{The Catholic Encyclopedia} (New York, 1908--1912): Gerard, ``General Chronology'', and Howlett, ``Biblical Chronology'' (vol.~3, 1908); Corbett, ``David'' (vol.~4, 1908); Souvay, ``Isaias'' (vol.~8, 1910); MacRory, ``Gospel of Saint Mark'' (vol.~9, 1910); Jacquier, ``Gospel of St.~Matthew'' (vol.~10, 1911); Prat, ``St.~Paul'' (vol.~11, 1911); Drum, ``Psalms'', and Vander Heeren, ``Epistle to the Philippians'' (vol.~12, 1911); Drum, ``Epistles to the Thessalonians'', Durand, ``The New Testament'', and Meistermann, ``Temple of Jerusalem'' (vol.~14, 1912); read in the New Advent text.
\item United States Conference of Catholic Bishops, \work{New American Bible, Revised Edition}, introductions to Isaiah and to the Psalms, \url{https://bible.usccb.org/bible/isaiah/0} and \url{https://bible.usccb.org/bible/psalms/0}, read in their public deliveries of 16~October 2026 (summarized; protected).
\item Pontifical Biblical Commission, \latin{De auctore, de tempore compositionis et de historica veritate Evangeliorum secundum Marcum et secundum Lucam}, \work{Acta Apostolicae Sedis} 4 (1912), pp.~463--465, responses I--IX, as the chronology corpus records them.
\end{itemize}
\par}
